\documentclass[11pt,a4paper]{article}

\usepackage[utf8]{inputenc}
\usepackage[T1]{fontenc}
\usepackage{lmodern}
\usepackage[ngerman]{babel}
\usepackage{microtype}
\usepackage{amsmath,amssymb,amsthm,mathtools}
\usepackage{geometry}
\usepackage{booktabs}
\usepackage{array}
\usepackage{enumitem}
\usepackage{xcolor}
\usepackage{csquotes}
\usepackage{setspace}
\usepackage[bottom]{footmisc}
\usepackage[colorlinks=true,linkcolor=blue!40!black,citecolor=blue!40!black,urlcolor=blue!40!black,
            pdftitle={Stern-Gerlach-Aufspaltung und Perihelpräzession in R3 und R8},
            pdfauthor={Thomas Hoffbauer}]{hyperref}

\geometry{a4paper,tmargin=2.5cm,bmargin=2.6cm,lmargin=2.4cm,rmargin=2.4cm}
\setstretch{1.08}
\setlist[itemize]{leftmargin=1.4em,itemsep=0.2em,topsep=0.3em}
\setlist[enumerate]{leftmargin=1.6em,itemsep=0.2em,topsep=0.3em}

\definecolor{historienblau}{RGB}{18,52,86}

\theoremstyle{definition}
\newtheorem{definition}{Definition}[section]
\theoremstyle{remark}
\newtheorem{bemerkung}{Bemerkung}[section]

\newcommand{\so}{\mathfrak{so}}
\newcommand{\gtwo}{\mathfrak{g}_2}
\newcommand{\ImO}{\operatorname{Im}}
\newcommand{\Spin}{\operatorname{Spin}}
\newcommand{\Aut}{\operatorname{Aut}}
\newcommand{\R}{\mathbb{R}}
\newcommand{\C}{\mathbb{C}}
\newcommand{\HH}{\mathbb{H}}
\newcommand{\OO}{\mathbb{O}}

\title{\textbf{Stern-Gerlach-Aufspaltung und Perihelpräzession}\\[0.3em]
\large in \(\R^3\) und \(\R^8\)}
\author{Thomas Hoffbauer}
\date{September 2026}

\begin{document}
\maketitle

\begin{abstract}
\noindent
Mit dem formalen Apparat der IABC-Erweiterung%
\footnote{In den übrigen Texten dieses Ordners erscheint dasselbe Rahmenmodell unter dem Kürzel
EABC (``Bamberger Modell''). IABC und EABC bezeichnen hier denselben Defekt-/Oktonionen-Apparat;
im vorliegenden Aufsatz wird die Notation des gegebenen Quelltextes (IABC) beibehalten.}
--- den Projektionen \(\gtwo\hookrightarrow\so(7)\) bzw.\ \(\so(4)\hookrightarrow\so(3)\) und der
feldtheoretischen Bivektorkopplung --- lässt sich die Übertragung auf das Stern-Gerlach-Experiment
sowie auf die Periheldrehung in \(\R^3\) und \(\R^8\) (bzw.\ dem 7D-Imaginärraum
\(\ImO(\OO)\cong\R^7\) und dem vollen oktonionischen Raum \(\OO\cong\R^8\)) exakt bestimmen.
Methodisch ist strikt zu trennen zwischen \emph{quantenmechanischer Spinor-Aufspaltung}
(Stern-Gerlach) und \emph{klassischer Geodäten- bzw.\ Bahnstörung} (Perihelpräzession).
Der Aufsatz reichert den gegebenen Apparat um historische Anker und kurze Herleitungen an und
fasst die Vorhersagen in einer Übersichtstabelle zusammen.
\end{abstract}

\tableofcontents
\vspace{1em}

%---------------------------------------------------------------------
\section{Einleitung: zwei Größen, zwei Ontologien}
%---------------------------------------------------------------------

Das Stern-Gerlach-Experiment und die Periheldrehung des Merkur gehören zu den klassischen
Prüfsteinen der modernen Physik --- der eine zur Raumquantisierung und später zum Spin,
der andere zur Geometrie der Gravitation. Beide lassen sich formal in denselben IABC-Apparat
einbetten, aber sie antworten auf unterschiedliche Fragen:

\begin{itemize}
  \item \textbf{Stern-Gerlach:} Wie zerlegt ein inhomogenes äußeres Feld einen internen
        Multipol- bzw.\ Spinor-Grad der Freiheit in räumlich getrennte Strahlen?
  \item \textbf{Perihel:} Wie verschiebt eine Abweichung vom Bertrandschen \(1/r\)-Potential
        die Apsidenlinie einer gebundenen klassischen Bahn?
\end{itemize}

Die erste Größe ist operatorwertig und diskret; die zweite ist eine säkulare Orbitmittelung
über eine klassische Phase. Eine Vermischung beider Lesarten --- etwa die Deutung einer
\(G_2\)-Verzweigung \(8\to 1\oplus 7\) als astronomische Merkurbeobachtung --- wäre
kategoriell falsch. Der vorliegende Text hält beide Stränge getrennt und vergleicht sie
nur am Ende in einer Tabelle.

%---------------------------------------------------------------------
\section{Historischer Rahmen}
%---------------------------------------------------------------------

\subsection{Stern-Gerlach, Spin und Dirac}

Im Februar 1922 führten Otto Stern und Walther Gerlach in Frankfurt das Experiment aus, das
heute ihren Namen trägt: ein Strahl von Silberatomen durchquerte ein stark inhomogenes
Magnetfeld und spaltete sich in \emph{zwei} diskrete Spuren auf dem Auffänger auf
\cite{SternGerlach1922}. Die zeitgenössische Deutung las dies als Raumquantisierung des
Bahnimpulses (\(l=1\)); die spätere Quantenmechanik erkannte darin den Elektronenspin
\(s=1/2\) des äußersten Elektrons im Silberggrundzustand (\(^2S_{1/2}\)), also \(2s+1=2\)
Komponenten.

Uhlenbeck und Goudsmit schlugen 1925 den Elektronenspin vor; Pauli formulierte die
zweikomponentigen Spinmatrizen \(\sigma_i\); Dirac vereinigte 1928 Spin und Relativität in
der vierkomponentigen Theorie. Der Stern-Gerlach-Effekt blieb das paradigmatische Beispiel
einer kraftartigen Aufspaltung durch
\begin{equation}
  \mathbf{F}=\nabla(\boldsymbol{\mu}\cdot\mathbf{B}).
\end{equation}

\subsection{Perihel, Einstein und Bertrand}

Urbain Le~Verrier hatte 1859 eine anomale Vorholung des Merkurperihels angezeigt. Albert
Einstein erklärte am 18.~November 1915 den Rest von rund \(43''\) pro Jahrhundert aus der
Allgemeinen Relativitätstheorie \cite{Einstein1915}. Die Schwarzschild-Metrik liefert das
effektive Störpotential \(\delta V_{\mathrm{GR}}\propto -1/r^3\) und damit die kanonische
Formel
\begin{equation}
  \Delta\varphi_{\mathrm{GR}}=\frac{6\pi GM}{c^2 a(1-e^2)}
  \label{eq:einstein}
\end{equation}
pro Umlauf. Bertrand zeigte 1873, dass unter allen zentralen Kräften in drei Raumdimensionen
nur der isotrope harmonische Oszillator und das Kepler-Potential \(1/r\) geschlossene
gebundene Bahnen erzeugen \cite{Bertrand1873}. Jede Abweichung --- relativistisch oder
anisotrop --- öffnet die Ellipse und erzeugt Präzession. Der Runge-Lenz-Vektor ist dann
nicht mehr konstant; seine säkulare Drehung ist die Periheldrift.

\subsection{Oktonionen, Hurwitz und \texorpdfstring{\(G_2\)}{G2}}

Die Oktonionen \(\OO\) wurden unabhängig von Graves (1843) und Cayley (1845) entdeckt.
Hurwitz' Satz besagt, dass die einzigen normierten Divisionsalgebren über \(\R\) die
Algebren \(\R,\C,\HH,\OO\) der Dimensionen \(1,2,4,8\) sind. Die Automorphismengruppe
\(G_2=\Aut(\OO)\) ist die Ausnahme-Lie-Gruppe vom Rang~2 und der Dimension~14; sie lässt
die reelle Einheit fest und wirkt irreduzibel auf dem imaginären Raum \(\ImO(\OO)\cong\R^7\).
Cartans Trialität verknüpft die drei achtdimensionalen Darstellungen von \(\Spin(8)\)
(\(\mathbf{8}_v\), \(\mathbf{8}_s\), \(\mathbf{8}_c\)) und ist eng mit der Multiplikation
in \(\OO\) verbunden \cite{Baez2002,Cartan1925}.

%---------------------------------------------------------------------
\section{I.\ Stern-Gerlach in \texorpdfstring{\(\R^3\)}{R\string^3}}
%---------------------------------------------------------------------

\subsection{Standardaufspaltung}

Im gewöhnlichen 3D-Laborraum wirkt der Hamilton-Operator aus dem IABC-Modell
\begin{equation}
  H=-\frac{\hbar^2}{2m}D_i D_i+V(x)+g_1 X^a\tau_a+H_Z,\qquad i\in\{1,2,3\}.
\end{equation}
Ein transversaler Gradient \(\partial_z B_z\) spaltet einen \(S=1/2\)-Strahl
(Silberatome, oder H-Atome im Grundzustand) entlang der Eigenwerte \(m_s=\pm 1/2\) von
\(\sigma_z\) in zwei Strahlen auf (\(2S+1=2\)).

\subsection{IABC-Korrektur: skalarer Drift, keine Spinaufspaltung erster Ordnung}

Der direkte Polarisationsterm \(g_1 X^a\tau_a\) koppelt an das interne Defektfeld. Besitzt
dieses einen räumlichen Gradienten (\(\partial_z X^a\neq 0\)), so entsteht eine zusätzliche
kraftartige Wirkung
\begin{equation}
  F_z^{\mathrm{IABC}}=g_1(\partial_z X^a)\,\tau_a.
  \label{eq:FzIABC}
\end{equation}

Warum gibt dies für \(J=1/2\) \emph{keine} Aufspaltung erster Ordnung? Die Operatoren
\(\tau_a\) wirken auf den betrachteten \(S\)-Zuständen als Vielfache der Identität (skalarer
Kanal): sie vertauschen mit den Paulimatrizen und besitzen in dem relevanten Unterraum nur
einen gemeinsamen Eigenwert. Folglich ist \(F_z^{\mathrm{IABC}}\) proportional zur
Einheitsmatrix im Spinraum. Beide \(m_s\)-Komponenten erfahren dieselbe Kraft und werden
gemeinsam ausgelenkt --- ein Schwerpunkt-Drift des ganzen Doppelstrahls, keine relative
Trennung.

Eine Phasenverschiebung im Interferometer entsteht, sobald entlang zweier räumlich
getrennter Pfade unterschiedliche Mittelwerte von \(X^a\) (oder unterschiedliche Verweilzeiten
im Gradientengebiet) akkumuliert werden: die Wirkung
\(\Delta S\sim g_1\int(X^a\tau_a)\,dt\) erzeugt eine relative Phase \(\Delta\phi=\Delta S/\hbar\),
auch wenn die Kraft in beiden Armen gleich und die Strahlaufspaltung null bleibt. Erst
Gradienten, die an \emph{nichtskalare} Operatoren (echte Spin-/Tensoroperatoren) koppeln,
liefern eine \(m\)-abhängige Kraftdifferenz.

%---------------------------------------------------------------------
\section{II.\ Stern-Gerlach in \texorpdfstring{\(\R^8\)}{R\string^8} / Oktonionen}
%---------------------------------------------------------------------

\subsection{Spinoren von \texorpdfstring{\(\Spin(8)\)}{Spin(8)} und Cartan-Trialität}

Wird das Experiment theoretisch in einem 8-dimensionalen Raum formuliert --- Bahndrehimpuls
in \(\so(8)\), interne Struktur klassifiziert durch \(G_2=\Aut(\OO)\) ---, so sind die
Spinoren von \(\Spin(8)\) achtdimensional: die beiden chiralen Halbmoduln \(\mathbf{8}_s\),
\(\mathbf{8}_c\) und die Vektordarstellung \(\mathbf{8}_v\), verknüpft über Cartans Trialität.

Ein inhomogenes verallgemeinertes Eichfeld \(N_\mu^A\) in \(\gtwo\) oder ein Gradient
\(\nabla_{\R^8}X^a\) mit sieben transversalen Richtungen erzeugt Kräfte entlang der
Cartan-Unteralgebra von \(\gtwo\) (Rang~2) bzw.\ von \(\so(8)\) (Rang~4):
\begin{equation}
  F_M=\nabla_{\R^8}\bigl(g_1 X^a\tau_a+N_\mu^A T_A\bigr).
\end{equation}

\subsection{Die Verzweigung \texorpdfstring{\(8\to 1\oplus 7\)}{8 -> 1 + 7} unter \texorpdfstring{\(G_2\)}{G2}}

Als reeller Vektorraum ist \(\OO=\R\cdot 1\oplus\ImO(\OO)\). Jeder Automorphismus
\(g\in G_2=\Aut(\OO)\) erfüllt \(g(1)=1\) und erhält die Multiplikation; insbesondere ist
die reelle Achse Fixgerade. Das Komplement \(\ImO(\OO)\) ist die fundamentale irreduzible
7-dimensionale Darstellung von \(G_2\). Damit zerfällt die 8-dimensionale
Oktonionen-Darstellung:
\begin{equation}
  \mathbf{8}\xrightarrow{G_2}\mathbf{1}\oplus\mathbf{7}.
\end{equation}

\subsection{Cartan-Ränge und die Zahl unabhängiger Gradienten}

\(\operatorname{Rang}(G_2)=2\) und \(\operatorname{Rang}(\so(8))=4\). Ein maximal
diagonalisierbares Stern-Gerlach-``Magnetfeld'' kann also höchstens so viele unabhängige
kommutierende Generatoren ansteuern, wie der Rang der betrachteten Lie-Algebra beträgt:
zwei für \(\gtwo\), vier für \(\so(8)\). Die oft zitierte Bildsprache ``sieben disjunkte
Strahlen plus ein Zentralstrahl'' ist eine \emph{darstellungstheoretische} Aussage, keine
empirische 8D-Beobachtung. Sie setzt voraus, dass der äußere Gradient die Darstellung
\(\mathbf{7}\) von \(\ImO(\OO)\) in Eigenräume zerlegt --- etwa durch Ausrichtung entlang
einer Cartan-Unteralgebra von \(\gtwo\) und eine geeignete Einbettung der sieben imaginären
Einheiten in Gewichtsräume. Der Zentralstrahl entspricht dem \(G_2\)-Singulett
\(\R\cdot 1\); die sieben Komponenten in \(\ImO(\OO)\) erfahren generisch nichttriviale
Kräfte. Ohne diese Ausrichtungsannahme bleibt nur die Verzweigung \(1\oplus 7\) als
Zerlegung, nicht notwendig eine siebenfache räumliche Aufspaltung.

%---------------------------------------------------------------------
\section{III.\ Perihel in \texorpdfstring{\(\R^3\)}{R\string^3}: ART und IABC-Quadrupol}
%---------------------------------------------------------------------

\subsection{Effektives Potential und Einstein-Präzession}

Das klassische Ein-Körper-Problem unter dem modifizierten IABC-Potential lautet
\begin{equation}
  V_{\mathrm{eff}}(r)=-\frac{k}{r}+\frac{L^2}{2m r^2}+\delta V_{G_2}(r,\theta),
\end{equation}
wobei \(\delta V_{G_2}\) aus dem quadratischen Defektterm \(g_2|X|^2\) und der effektiven
Rang-2-Torsion
\begin{equation}
  H_Q=\frac{\varepsilon_A}{2m}Q_{ij}p_i p_j
\end{equation}
stammt. Die Schwarzschild-Metrik liefert das effektive Störpotential
\begin{equation}
  \delta V_{\mathrm{GR}}(r)=-\frac{GM L^2}{c^2 m r^3},
\end{equation}
was zur Einstein-Präzession~\eqref{eq:einstein} führt.

\emph{Skizze zum Faktor \(6\pi\).}
 In der säkularen Störungstheorie für eine kleine zentrale
Korrektur \(\delta V\sim -\alpha/r^3\) (hier \(\alpha=GM L^2/(c^2 m)\)) dreht sich der
Runge-Lenz-Vektor mit einer Rate proportional zu \(\partial(\delta V)/\partial r\) bzw.\ zum
Orbitmittel von \(r^{-3}\). Für eine Kepler-Ellipse mit Halbachse \(a\) und Exzentrizität
\(e\) gilt \(\langle r^{-3}\rangle=1/\bigl[a^3(1-e^2)^{3/2}\bigr]\). Zusammen mit dem
Faktor aus der Apsidendrehung (zwei Radialumkehrpunkte pro Umlauf, Normierung auf \(2\pi\))
ergibt sich das Standardresultat \(\Delta\varphi=6\pi GM/\bigl(c^2 a(1-e^2)\bigr)\). Die
genaue Buchhaltung steht in den Lehrbüchern der ART \cite{Weinberg1972,MTW1973}; hier wird
keine abweichende Herleitung behauptet.

\subsection{IABC-\texorpdfstring{\(G_2\)}{G2}-Torsionsbeitrag}

Die Kopplung an das Rang-2-Vakuumfeld \(Q_{ij}=\operatorname{diag}(1,-1,0)\) bricht die
sphärische Symmetrie und induziert eine azimutale Störung
\begin{equation}
  \delta V_Q(r,\varphi)=\frac{\varepsilon_A}{2}\,\frac{L^2}{m r^2}\cos(2\varphi).
\end{equation}
Über einen geschlossenen Kepler-Orbit ist der Mittelwert von \(\cos(2\varphi)\)
\emph{nicht} automatisch eine Periheldrehung derselben Art wie die ART-Korrektur
\(\propto r^{-3}\). Die quadrupolare Anisotropie erzeugt vielmehr eine säkulare Drift, deren
Amplitude und Phasenlage vom Anfangsperihel \(\varphi_0\) und von Orbitmitteln der Form
\(\langle\cos(2\varphi)\rangle\) abhängen. Im gegebenen IABC-Apparat lautet die
Modellaussage
\begin{equation}
  \Delta\varphi_{G_2}^{(3)}\propto\varepsilon_A\,C_2\cos(2\varphi_0),
  \label{eq:G2drift}
\end{equation}
mit einem dimensionslosen Orbitfaktor \(C_2\). Gleichung~\eqref{eq:G2drift} ist eine
Vorhersage des Modells, keine gemessene Abweichung von der Einstein-Präzession.

Astronomische Tests der ART (Merkur, Venus, Radarlaufzeiten) bestätigen die
Einstein-Formel typischerweise auf der Größenordnung relativer Residuen von
\(\sim 10^{-4}\) (je nach Epoche und Analyse). Daraus folgt qualitativ eine makroskopische
Schranke an den Tensorparameter \(\varepsilon_A\): \(|\varepsilon_A|\) muss so klein sein,
dass der Modellbeitrag~\eqref{eq:G2drift} innerhalb der üblichen Merkur-Residuen bleibt.
Eine scharfe numerische Obergrenze für \(\varepsilon_A\) wird hier nicht neu hergeleitet;
die Aussage bleibt eine Abschätzung der Form
\begin{equation}
  |\varepsilon_A|\ll\frac{\text{(übliche Merkur-Residuen)}}{|C_2\cos(2\varphi_0)|},
\end{equation}
solange \(C_2\cos(2\varphi_0)\) nicht pathologisch klein ist.

%---------------------------------------------------------------------
\section{IV.\ Perihel in \texorpdfstring{\(\R^8\)}{R\string^8}: kein Kepler, Defektpotential}
%---------------------------------------------------------------------

\subsection{Gaußsches Gesetz und Instabilität}

In Dimensionen \(D\geq 4\) ergibt das Gaußsche Gesetz im euklidischen Raum \(\R^D\) für
Punktquellen
\begin{equation}
  F(r)\propto -\,r^{1-D}\qquad\Rightarrow\qquad
  V(r)\propto
  \begin{cases}
    r^{2-D}, & D\neq 2,\\
    \log r, & D=2.
  \end{cases}
\end{equation}
Für \(D=8\) ist das ungestörte Zentralpotential \(V(r)=-k_8/r^6\). Mit Zentrifugalbarriere
lautet das effektive Radialpotential
\begin{equation}
  V_{\mathrm{eff}}(r)=-\frac{k_8}{r^6}+\frac{L^2}{2m r^2}.
  \label{eq:Veff8}
\end{equation}

\emph{Warum kein stabiles gebundenes Minimum?} Betrachte
\begin{equation}
  V_{\mathrm{eff}}'(r)=\frac{6k_8}{r^7}-\frac{L^2}{m r^3}.
\end{equation}
Eine kritische Stelle \(V_{\mathrm{eff}}'(r_*)=0\) existiert formal bei
\(r_*^4=6k_8 m/L^2\) (für \(k_8,L\neq 0\)). Die zweite Ableitung
\begin{equation}
  V_{\mathrm{eff}}''(r)=\frac{L^2}{m r^4}-\frac{42 k_8}{r^8}
\end{equation}
ergibt bei \(r_*\) den Wert \(V_{\mathrm{eff}}''(r_*)=-6L^2/(m r_*^4)<0\): der kritische Punkt
ist ein \emph{Maximum}, kein Minimum. Äquivalent: für \(r\to 0\) dominiert \(-1/r^6\) die
Zentrifugalbarriere \(+1/r^2\), und \(V_{\mathrm{eff}}\to-\infty\); für \(r\to\infty\) gilt
\(V_{\mathrm{eff}}\to 0^-\) bzw.\ von oben, je nach Vorzeichenkonvention der Energie.
Gebundene Ellipsen mit zwei Umkehrpunkten gibt es damit im reinen Problem~\eqref{eq:Veff8}
nicht: Teilchen stürzen spiralförmig ins Zentrum oder entkommen ins Unendliche.

\subsection{Kompaktifizierung und IABC-Defektpotential}

Damit in acht Dimensionen geschlossene oder quasigeschlossene Bahnen existieren, über die
eine Periheldrehung überhaupt definiert werden kann, darf der Raum nicht aus acht
unkompaktifizierten flachen Dimensionen bestehen. Das IABC-Schema nutzt Zerlegungen der Art
\begin{equation}
  \R^8\cong\R^4_{\mathrm{Raumzeit}}\times\HH^\perp
  \qquad\text{bzw.}\qquad
  \R^3_{\mathrm{Raum}}\times S^1_{\mathrm{Zeit}}\times M_{\mathrm{int}}^4.
\end{equation}
Wird die 8-dimensionale Bewegung auf eine 3-Sphäre oder durch den \(G_2\)-Hedgehog
topologisch gefangen, wirkt das Defektpotential
\begin{equation}
  V_{\mathrm{IABC}}(r)=V_{\mathrm{classical}}(r)+g_2|X(r)|^2.
\end{equation}
Mit einem Kink-Profil \(|X(r)|^2=v^2\bigl(1-\cdots\bigr)\) entsteht eine rücktreibende Kraft,
die Umkehrpunkte \(r_{\min}\), \(r_{\max}\) erzeugen kann.

Die Perihelverschiebung berechnet sich dann über das Bahnintegral (Runge-Lenz-/Apsidenform)
\begin{equation}
  \Delta\varphi^{(8)}
  =2\int_{r_{\min}}^{r_{\max}}
  \frac{L/(m r^2)\,dr}{\sqrt{(2/m)\bigl(E-V_{\mathrm{eff}}^{(8)}(r)\bigr)}}
  -2\pi.
  \label{eq:Dphi8}
\end{equation}
In unkompaktifiziertem \(\R^8\) mit reinem \(-1/r^6\)-Potential fehlen die Umkehrpunkte:
die Wurzel unter dem Integral wird nicht auf einem kompakten Intervall reell und
positiv. Erst das zusätzliche Defektpotential \(g_2|X(r)|^2\) (oder eine andere
effektive Bindung durch Kompaktifizierung) macht \(E-V_{\mathrm{eff}}^{(8)}\) auf einem
Intervall \([r_{\min},r_{\max}]\) nichtnegativ und damit \(\Delta\varphi^{(8)}\)
wohldefiniert.

%---------------------------------------------------------------------
\section{Fazit}
%---------------------------------------------------------------------

\begin{center}
\small
\renewcommand{\arraystretch}{1.3}
\begin{tabular}{@{}>{\raggedright\arraybackslash}p{0.20\textwidth}
                >{\raggedright\arraybackslash}p{0.36\textwidth}
                >{\raggedright\arraybackslash}p{0.36\textwidth}@{}}
\toprule
\textbf{Größe} & \textbf{in \(\R^3\)} & \textbf{in \(\R^8\)} \\
\midrule
Stern-Gerlach &
Zweifach-Aufspaltung \((\pm\hbar/2)\); IABC: skalarer Drift bzw.\ Phasenverschiebung über \(g_1 X^a\tau_a\) &
Unter \(G_2\subset\mathrm{SO}(8)\): \(\mathbf{1}\oplus\mathbf{7}\) (Zentralstrahl + 7 Komponenten auf \(\ImO(\OO)\); darstellungstheoretisch) \\
\addlinespace
Perihel\-präzession &
Einstein-Präzession plus quadrupolare \(\cos(2\varphi)\)-Anisotropie aus \(Q_{ij}\) (Modell: \(\propto\varepsilon_A C_2\cos(2\varphi_0)\)) &
Reine \(\R^8\)-Keplerbahnen kollabieren (\(V\sim -1/r^6\)). Präzession nur im kompakten/solitären IABC-Defektpotential \\
\bottomrule
\end{tabular}
\end{center}

\medskip
Der mathematische Apparat liefert für beide Größen wohldefinierte Vorhersagen. In \(\R^3\)
fungieren sie als experimentelle Schranken für Laborspektroskopie und Himmelsmechanik; in
\(\R^8\) beschreiben sie die geometrische Entfaltung der Oktonionen-Algebra in ihre
fundamentalen Darstellungen --- nicht die Existenz eines acht-dimensionalen
Laborstrahls oder einer acht-dimensionalen Merkurbahn im flachen Raum.

%---------------------------------------------------------------------
\begin{thebibliography}{99}

\bibitem{SternGerlach1922}
O.~Stern und W.~Gerlach,
\emph{Der experimentelle Nachweis der Richtungsquantelung im Magnetfeld},
Zeitschrift für Physik \textbf{9} (1922), 349--352.

\bibitem{Einstein1915}
A.~Einstein,
\emph{Erklärung der Perihelbewegung des Merkur aus der allgemeinen Relativitätstheorie},
Sitzungsberichte der Königlich Preußischen Akademie der Wissenschaften (Berlin) (1915),
831--839.

\bibitem{Bertrand1873}
J.~Bertrand,
\emph{Théorème relatif au mouvement d'un point attiré vers un centre fixe},
Comptes Rendus de l'Académie des Sciences \textbf{77} (1873), 849--853.

\bibitem{Hurwitz1898}
A.~Hurwitz,
\emph{Über die Composition der quadratischen Formen von beliebig vielen Variablen},
Nachrichten von der Gesellschaft der Wissenschaften zu Göttingen, Mathematisch-Physikalische
Klasse (1898), 309--316.

\bibitem{Cartan1925}
É.~Cartan,
\emph{Le principe de dualité et la théorie des groupes simples et semi-simples},
Bulletin des Sciences Mathématiques \textbf{49} (1925), 361--374.
% Trialität und Spin(8): klassische Cartan-Arbeiten; vgl. auch die Darstellung in Baez.

\bibitem{Baez2002}
J.~C.~Baez,
\emph{The Octonions},
Bulletin of the American Mathematical Society \textbf{39} (2002), 145--205.

\bibitem{Weinberg1972}
S.~Weinberg,
\emph{Gravitation and Cosmology: Principles and Applications of the General Theory of
Relativity},
Wiley, New York, 1972.

\bibitem{MTW1973}
C.~W.~Misner, K.~S.~Thorne und J.~A.~Wheeler,
\emph{Gravitation},
W.~H.~Freeman, San Francisco, 1973.

\end{thebibliography}

\end{document}
